\documentclass[11pt,a4paper]{article}

% ---------------------------------------------------------
% PREAMBLE (stable, warning-free, Overleaf-safe)
% ---------------------------------------------------------
\usepackage[utf8]{inputenc}
\usepackage[T1]{fontenc}
\usepackage{lmodern}
\usepackage{amsmath, amssymb, amsfonts}
\usepackage{bm}
\usepackage{geometry}
\usepackage{graphicx}
\usepackage{hyperref}
\usepackage{cite}
\usepackage{authblk}
\usepackage{physics}
\usepackage{mathtools}

\geometry{margin=2.4cm}

\title{\textbf{Companion LXIII: Persistent Homology of Weak-Lensing Fields in the Relaxation-Driven Cyclic Cosmology}}
\author{Michael Lehmann \\ \texttt{mi.lehmann@gmx.de}}
\date{April 2026}

\begin{document}
\maketitle

\begin{abstract}
Persistent homology provides a topological characterization of cosmological fields 
by tracking the birth and death of connected components, loops, and void-like 
structures across a continuous filtration. In weak-lensing analyses, persistence 
diagrams and barcodes quantify the stability of topological features in the 
convergence field. In the standard cosmological model, persistence reflects the 
nonlinear matter distribution and the projection of large-scale structure. In the 
Relaxation-Driven Cyclic Cosmology (RDCC), the scale-dependent evolution of the 
gravitational potential modifies the stability, lifetime, and hierarchical merging 
of topological features in a characteristic way. This Companion develops a 
continuous description of persistent homology in RDCC, deriving how the relaxation 
scale influences persistence diagrams, barcode statistics, and the observational 
signatures in weak-lensing surveys. The resulting framework provides a distinctive 
topological probe of the RDCC gravitational field.
\end{abstract}
% ---------------------------------------------------------
\section{Introduction}

Weak-lensing convergence maps contain rich topological information. While Betti 
numbers count topological features at fixed thresholds, persistent homology tracks 
their evolution across a continuous filtration. This provides a powerful probe of 
non-Gaussianity, hierarchical structure, and the stability of topological features.

In the standard cosmological model, persistence diagrams reflect the nonlinear 
matter distribution and the projection of large-scale structure. In RDCC, the 
gravitational potential evolves differently. The relaxation mechanism suppresses 
the decay of small-scale modes and enhances the coherence of large-scale modes. 
This modifies the birth, death, and lifetime of topological features in a 
scale-dependent manner, producing a distinctive signature in persistent-homology 
analyses.
% ---------------------------------------------------------
\section{Persistent Homology of the Convergence Field}

For a 2D field such as the convergence $\kappa(\hat{\mathbf{n}})$, persistent 
homology tracks the evolution of topological features across a filtration 
parameter $\nu$:

\begin{itemize}
    \item $\beta_{0}$ features (connected components) are born when isolated 
          high-$\kappa$ regions appear.
    \item $\beta_{1}$ features (loops) are born when ring-like structures form.
    \item Features die when they merge into larger structures or collapse.
\end{itemize}

The persistence diagram consists of points $(\nu_{\mathrm{birth}},\nu_{\mathrm{death}})$.
% ---------------------------------------------------------
\section{RDCC Modification of the Convergence Field}

In RDCC, the convergence evolves as
\begin{equation}
    \kappa_{\mathrm{RDCC}}(\ell)
    = \kappa(\ell)
      \left[
        1 - \chi_{\kappa}
        \left(\frac{\ell}{\ell_{\mathrm{rel}}}\right)^{\alpha}
      \right].
\end{equation}

The variance becomes
\begin{equation}
    \sigma_{\kappa,\mathrm{RDCC}}^{2}
    = \sigma_{\kappa}^{2}
      \left[
        1 - \chi_{\sigma}
        \left(\frac{\ell}{\ell_{\mathrm{rel}}}\right)^{\alpha}
      \right].
\end{equation}

These modifications alter the birth and death thresholds of topological features.
% ---------------------------------------------------------
\section{Persistent Homology in RDCC}

The RDCC persistence diagram becomes
\begin{equation}
    D_{\mathrm{RDCC}}
    = \left\{
        \left(
          \nu_{\mathrm{birth}},
          \nu_{\mathrm{death}}
        \right)
        \left[
          1 - \chi_{D}
          \left(\frac{\ell}{\ell_{\mathrm{rel}}}\right)^{\alpha}
        \right]
      \right\}.
\end{equation}

The lifetime of a feature is
\begin{equation}
    L = \nu_{\mathrm{death}} - \nu_{\mathrm{birth}},
\end{equation}

and in RDCC,
\begin{equation}
    L_{\mathrm{RDCC}}
    = L
      \left[
        1 - \chi_{L}
        \left(\frac{\ell}{\ell_{\mathrm{rel}}}\right)^{\alpha}
      \right].
\end{equation}

This modification reflects the relaxation-driven changes in topological stability.
% ---------------------------------------------------------
\section{Topological Signatures of RDCC}

The relaxation mechanism enhances the coherence of large-scale modes. Persistent 
homology therefore exhibits:

\begin{itemize}
    \item longer-lived large-scale connected components,
    \item suppressed short-lived small-scale features,
    \item fewer but more persistent loops,
    \item a shift in the persistence-lifetime distribution,
    \item a characteristic turnover near $\ell_{\mathrm{rel}}$,
    \item modified hierarchical merging of topological structures.
\end{itemize}

These signatures provide a direct observational probe of the RDCC gravitational 
field in persistence space.
% ---------------------------------------------------------
\section{Conclusion}

Persistent homology of weak-lensing fields in the Relaxation-Driven Cyclic 
Cosmology reflects the scale-dependent evolution of the gravitational potential and 
the nonlinear matter distribution. The relaxation mechanism modifies the birth, 
death, and lifetime of topological features in a scale-dependent manner, producing 
distinctive signatures in weak-lensing surveys. These effects provide a direct 
observational window into the relaxation scale and the temporal structure of the 
RDCC gravitational field. The present Companion establishes the theoretical 
foundation for future studies of topological cosmology, nonlinear structure, and 
the interplay between light and gravity in RDCC.

\bibliographystyle{apsrev4-2}
\nocite{*}
\bibliography{references}


\end{document}
